\section{The Quarrel over Novelty}

The Hebrew project was opposed in Jerome's lifetime by serious people with
serious arguments, and the opposition is documented on both sides of two
famous correspondences.

\textbf{Augustine's case.} Writing from Hippo in 403, Augustine thanked
Jerome warmly for the Gospel revision---``in almost every passage we have
found nothing to object to, when we compared it with the Greek
Scriptures''---and in the same letter begged him not to put the Hebrew
Bible into the churches. His grounds were pastoral and ecumenical: the
Septuagint bound Greek and Latin Christendom together, and a Latin text
answerable only to the Hebrew would make the Greek churches, and the
Jews, the only referees. His exhibit was a riot. A brother bishop---``it
was in the town of Oea,'' modern Tripoli---had introduced Jerome's
version; when the reading reached Jonah's plant---``a very different
rendering from that which had been of old familiar to the senses and
memory of all the worshippers''---the congregation erupted, ``especially
among the Greeks''; the local Jews, consulted, said (``whether from
ignorance or from spite,'' Augustine allows) that the Hebrew supported the
old reading; and the bishop ``was compelled to correct your version in
that passage as if it had been falsely translated, as he desired not to be
left without a congregation.''\endnote{Augustine, \work{Letter} 71.5--6
(A.D.~403), tr.\ Cunningham (NPNF series~1, vol.~1), checked at New
Advent. The request for a Latin translation of the Septuagint instead is
71.6. The anecdote is Augustine's report of a report; the offending word
(gourd versus ivy, Jonah 4:6) is identified in Jerome's reply.}

\textbf{Jerome's defense.} The reply (his Letter 112) concedes nothing.
The old dispute, he says, is ``about the gourd'': where the Hebrew has
\textit{ciceion} (the castor plant), the Seventy rendered ``gourd''
(\textit{cucurbita}); Aquila and the others, ``ivy'' (\textit{hedera});
transliterating would have meant nothing to Latin readers, gourd ``is not
found in the Hebrew,'' so he wrote ivy ``that I might not differ from all
other translators.'' On the larger charge he states his program in a
sentence: ``I have laboured not to supersede what has been long esteemed,
but only to bring prominently forward those things which have been either
omitted or tampered with by the Jews, in order that Latin readers might
know what is found in the original Hebrew''---and if anyone dislikes it,
``none compels him against his will'' to read it.\endnote{Jerome,
\work{Letter} 112.20--22 (= Augustine, \work{Letter} 75), tr.\
Cunningham (NPNF series~1, vol.~1), checked at New Advent. The
``tampered with by the Jews'' framing is Jerome's polemic, not this
account's judgment; his own prefaces elsewhere ground differences in
text and idiom rather than fraud.} The received Vulgate text of Jonah
4:6 duly reads \textit{hederam}---the ivy the Oea congregation
rioted over---and its standard printed editions carry it to this
day.\endnote{Checked visually at Jonah 4:6 in the Hetzenauer 1914
printing of the Clementine Vulgate (``Et pr\ae paravit Dominus Deus
hederam\ldots''), page 875 of the SacredBible page-image edition; the
1914 editor's marginal apparatus glosses the plant as \textit{ricinus},
quietly agreeing with Jerome's Hebrew botany while printing his ivy. See
the source records bound to this publication.}

\textbf{The settlement.} Augustine's last long letter in the exchange
(405) accepts the Hebrew translation's usefulness for scholarship while
restating the pastoral line: he asks for Jerome's Septuagint-based
revision precisely so that churches need not risk ``bringing forward
anything which was, as it were, new and opposed to the authority of the
Septuagint version.'' Fifteen years later, in the \work{City of God}, he
found the durable formula: Jerome is ``a man most learned, and skilled in
all three languages, who translated these same Scriptures into the Latin
speech, not from the Greek, but from the Hebrew''---a labor whose
faithfulness, Augustine notes, even the Jews acknowledge---and yet the
Church rightly holds to the Seventy, whose divergences the one Spirit may
have intended.
Both things, for Augustine, were true at once.\endnote{Augustine,
\work{Letter} 82.34--35, tr.\ Cunningham (NPNF series~1, vol.~1), checked
at New Advent; \work{City of God} 18.43, Marcus Dods translation
(Edinburgh: T.~\&~T. Clark, 1871), quoted wording verified against the
exact 1871 text registered in this repository's source library, and
present likewise in the NPNF series~1, vol.~2 reprint at CCEL. Augustine's
inspiration theology of the Seventy is his own argument, labeled as
such.}

\textbf{Rufinus's attack.} Where Augustine argued, Rufinus of
Aquileia---Jerome's estranged former friend---prosecuted. His
\work{Apology} calls the Hebrew project ``havoc'' wrought on ``the divine
record handed down to the Churches by the Apostles'': Susanna ``cut out,
thrown aside and dismissed,'' the Hymn of the Three Children ``erased
from the place where it stood''; the seventy translators in their cells
wrote by the Holy Spirit, but this is ``a translation made by a single
man under the inspiration of Barabbas''---and did Peter, who governed
the Roman church and delivered it its Scriptures, deceive it for four
hundred years? Now, he sneers, ``let us change the inscriptions upon the
tombs of the ancients, so that it may be known\ldots\ that it was not a
gourd but an ivy plant under whose shade Jonah
rested.''\endnote{Rufinus, \work{Apology against Jerome} 2.33--35, tr.\
Fremantle (NPNF series~2, vol.~3), checked at New Advent. The
``Barabbas'' gibe twists Jerome's own report that he had a Jewish
teacher, Baraninas (\work{Letter} 84.3). Rufinus writes as a party to a
bitter feud over Origen; his testimony is polemic, and is used here as
evidence of the opposition's arguments, not of Jerome's practice.}

The pattern of the resistance deserves notice: nobody defended the Old
Latin as a translation. The defenders of the old text defended the
authority of the Septuagint, the unity of Greek and Latin Christendom,
and the devotional memory of congregations---reception arguments, not
philological ones. That is why the quarrel could end, over centuries, in
a quiet reversal: the Church took Jerome's Bible, and kept the
Septuagint's canon.\endnote{Class~S synthesis over the letters cited in
this section.}
